\documentclass{article}
\usepackage[utf8]{inputenc}
\usepackage{amsmath,amssymb}
\usepackage[most]{tcolorbox}
\usepackage{geometry}
\geometry{margin=2.5cm}

\newcommand{\step}[1]{\par\noindent\hangindent=3.5em\hangafter=1%
  \makebox[3.5em][l]{\textbf{#1.}}}
\newcommand{\steptag}[1]{\hfill{\small\textsf{[#1]}}}

\newtcolorbox{proofbox}[1]{
  colback=gray!5, colframe=gray!60,
  fonttitle=\bfseries, title={#1},
  breakable, enhanced,
  left=4pt, right=4pt, top=4pt, bottom=4pt,
  before skip=10pt, after skip=10pt
}

\begin{document}

\begin{proofbox}{Witness arithmetic for PRH sharpness: for every real lamb...}
\step{1} Witness arithmetic for PRH sharpness: for every real lambda with 0 < lambda < 1/2, let A\_lambda in R\textasciicircum{}(4x2) have rows (1,0), (0,1), (1-lambda,lambda), (lambda,1-lambda), and let M\_lambda in R\textasciicircum{}(2x4) have rows (1-lambda,0,lambda,0), (0,1-lambda,0,lambda); A\_lambda and M\_lambda have nonnegative entries and probability-vector rows and hence represent positive unital maps A\_lambda:l-infinity(2)->l-infinity(4) and M\_lambda:l-infinity(4)->l-infinity(2); for every pair of integers r,s >= 1 and every B in R\textasciicircum{}(r x s), the induced l-infinity operator norm is ||B|$|\_\{infinity-\rangle$infinity\}=max\_\{1<=i<=r\} sum\_\{j=1\}\textasciicircum{}s |B\_ij|; M\_lambda A\_lambda has rows (1-lambda\textasciicircum{}2,lambda\textasciicircum{}2), (lambda\textasciicircum{}2,1-lambda\textasciicircum{}2), so epsilon\_lambda:=||M\_lambda A\_lambda-I\_2|$|\_\{infinity-\rangle$infinity\}=2*lambda\textasciicircum{}2 and epsilon\_lambda tends to 0 as lambda tends to 0; and P\_lambda:=A\_lambda M\_lambda is row-stochastic with rows (1-lambda,0,lambda,0), (0,1-lambda,0,lambda), ((1-lambda)\textasciicircum{}2,lambda*(1-lambda),lambda*(1-lambda),lambda\textasciicircum{}2), (lambda*(1-lambda),(1-lambda)\textasciicircum{}2,lambda\textasciicircum{}2,lambda*(1-lambda)), and ||row\_1(P\_lambda)-row\_3(P\_lambda)||\_1=2*lambda. \steptag{validated} \\[6pt]
\step{1.1} Fix an arbitrary real lambda with 0 < lambda < 1/2. Then lambda > 0 and 1-lambda > 1/2 > 0, so lambda, 1-lambda, lambda\textasciicircum{}2, (1-lambda)\textasciicircum{}2, and lambda(1-lambda) are all nonnegative. \steptag{validated} \\[4pt]
\step{1.2} For A\_lambda with rows (1,0), (0,1), (1-lambda,lambda), (lambda,1-lambda), all entries are nonnegative by 1.1, and the four row sums are respectively 1, 1, (1-lambda)+lambda=1, and lambda+(1-lambda)=1; hence every row of A\_lambda is a probability vector. \steptag{validated} \\[4pt]
\step{1.3} For M\_lambda with rows (1-lambda,0,lambda,0) and (0,1-lambda,0,lambda), all entries are nonnegative by 1.1, and each row sum is (1-lambda)+lambda=1; hence every row of M\_lambda is a probability vector. \steptag{validated} \\[4pt]
\step{1.4} By the probability-row equivalence in def-positive-approximate-retract, 1.2 and 1.3 imply that A\_lambda and M\_lambda represent positive unital maps A\_lambda:l-infinity(2)->l-infinity(4) and M\_lambda:l-infinity(4)->l-infinity(2), respectively. \steptag{validated} \\[4pt]
\step{1.5} For every pair of integers r,s >= 1 and every real r-by-s matrix B, its induced l-infinity operator norm satisfies ||B|$|\_\{infinity-\rangle$infinity\}=max\_\{1<=i<=r\} sum\_\{j=1\}\textasciicircum{}s |B\_ij|. \steptag{validated} \\[6pt]
\step{1.5.1} Let R:=max\_\{1<=i<=r\} sum\_\{j=1\}\textasciicircum{}s |B\_ij|. For every x in R\textasciicircum{}s with ||x||\_infinity<=1 and every i, |(Bx)\_i|=|sum\_j B\_ij x\_j|<=sum\_j |B\_ij||x\_j|<=sum\_j |B\_ij|<=R. Hence ||Bx||\_infinity<=R, and taking the supremum in the definition of the induced norm gives ||B|$|\_\{infinity-\rangle$infinity\}<=R. \steptag{validated} \\[4pt]
\step{1.5.2} Because r>=1 and the set of rows is finite, choose i\_0 attaining R. Define x\_j=1 if B\_(i\_0 j)>=0 and x\_j=-1 if B\_(i\_0 j)<0. Then ||x||\_infinity=1 and (Bx)\_(i\_0)=sum\_j |B\_(i\_0 j)|=R, so ||B|$|\_\{infinity-\rangle$infinity\}>=||Bx|$|\_infinity\rangle$=R. Together with 1.5.1 this proves the equality in 1.5. \steptag{validated} \\[4pt]
\step{1.5.3} For avoidance of ambiguity, put R:=max\_\{1<=i<=r\} sum\_\{j=1\}\textasciicircum{}s |B\_ij|. If ||x||\_infinity<=1, then |x\_j|<=1 for every j, and for each i the triangle inequality gives |(Bx)\_i|=|sum\_j B\_ij x\_j|<=sum\_j |B\_ij||x\_j|<=sum\_j |B\_ij|<=R. Therefore ||Bx||\_infinity<=R for every x with ||x||\_infinity<=1, so the definition of the induced norm yields ||B|$|\_\{infinity-\rangle$infinity\}<=R. In the challenge test B=[1], x=(1), R=1, this reads 1<=1<=1<=1; the challenge arose by dropping the equals signs from the <= symbols already present in node 1.5.1. \steptag{validated} \\[4pt]
\step{1.5.4} Put R:=max\_\{1<=i<=r\} sum\_\{j=1\}\textasciicircum{}s |B\_ij|. For any x in R\textasciicircum{}s with ||x||\_infinity<=1, each |x\_j|<=1, so for every row i the triangle inequality gives |(Bx)\_i|=|sum\_j B\_ij x\_j|<=sum\_j |B\_ij||x\_j|<=sum\_j |B\_ij|<=R. Taking the maximum over i and then the supremum over the l-infinity unit ball yields ||B|$|\_\{infinity-\rangle$infinity\}<=R. All comparisons are non-strict (for B=[1] and x=(1), each is equality). Conversely, because r>=1 is finite, choose i\_0 attaining R and set x\_j=1 when B\_(i\_0 j)>=0 and x\_j=-1 otherwise. Since s>=1, ||x||\_infinity=1, while (Bx)\_(i\_0)=sum\_j |B\_(i\_0 j)|=R; hence ||B|$|\_\{infinity-\rangle$infinity\}>=||Bx|$|\_infinity\rangle$=R. Thus ||B|$|\_\{infinity-\rangle$infinity\}=R, including the case R=0. \steptag{validated} \\[4pt]
\step{1.6} Direct row-by-column multiplication gives M\_lambda A\_lambda with first row (1-lambda)(1,0)+lambda(1-lambda,lambda)=(1-lambda\textasciicircum{}2,lambda\textasciicircum{}2) and second row (1-lambda)(0,1)+lambda(lambda,1-lambda)=(lambda\textasciicircum{}2,1-lambda\textasciicircum{}2). \steptag{validated} \\[6pt]
\step{1.6.1} The first row of M\_lambda is (1-lambda,0,lambda,0), so the first row of M\_lambda A\_lambda is (1-lambda)row\_1(A\_lambda)+lambda row\_3(A\_lambda)=(1-lambda)(1,0)+lambda(1-lambda,lambda)=(1-lambda+lambda-lambda\textasciicircum{}2,lambda\textasciicircum{}2)=(1-lambda\textasciicircum{}2,lambda\textasciicircum{}2). \steptag{validated} \\[4pt]
\step{1.6.2} The second row of M\_lambda is (0,1-lambda,0,lambda), so the second row of M\_lambda A\_lambda is (1-lambda)row\_2(A\_lambda)+lambda row\_4(A\_lambda)=(1-lambda)(0,1)+lambda(lambda,1-lambda)=(lambda\textasciicircum{}2,1-lambda+lambda-lambda\textasciicircum{}2)=(lambda\textasciicircum{}2,1-lambda\textasciicircum{}2). These are all rows of the 2-by-2 product, proving 1.6. \steptag{validated} \\[4pt]
\step{1.7} By 1.6, M\_lambda A\_lambda-I\_2 has rows (-lambda\textasciicircum{}2,lambda\textasciicircum{}2) and (lambda\textasciicircum{}2,-lambda\textasciicircum{}2); by the max-row formula 1.5, epsilon\_lambda:=||M\_lambda A\_lambda-I\_2|$|\_\{infinity-\rangle$infinity\}=max(2lambda\textasciicircum{}2,2lambda\textasciicircum{}2)=2lambda\textasciicircum{}2. \steptag{validated} \\[4pt]
\step{1.8} Since epsilon\_lambda=2lambda\textasciicircum{}2 by 1.7 and the polynomial 2lambda\textasciicircum{}2 tends to 0 as real lambda tends to 0, epsilon\_lambda tends to 0 as lambda tends to 0 (in particular along 0<lambda<1/2). \steptag{validated} \\[4pt]
\step{1.9} Direct row-by-matrix multiplication gives P\_lambda:=A\_lambda M\_lambda with rows (1-lambda,0,lambda,0), (0,1-lambda,0,lambda), ((1-lambda)\textasciicircum{}2,lambda(1-lambda),lambda(1-lambda),lambda\textasciicircum{}2), and (lambda(1-lambda),(1-lambda)\textasciicircum{}2,lambda\textasciicircum{}2,lambda(1-lambda)). \steptag{validated} \\[6pt]
\step{1.9.1} Because row\_1(A\_lambda)=(1,0), row\_1(A\_lambda M\_lambda)=row\_1(M\_lambda)=(1-lambda,0,lambda,0). \steptag{validated} \\[4pt]
\step{1.9.2} Because row\_2(A\_lambda)=(0,1), row\_2(A\_lambda M\_lambda)=row\_2(M\_lambda)=(0,1-lambda,0,lambda). \steptag{validated} \\[4pt]
\step{1.9.3} Because row\_3(A\_lambda)=(1-lambda,lambda), row\_3(A\_lambda M\_lambda)=(1-lambda)row\_1(M\_lambda)+lambda row\_2(M\_lambda)=((1-lambda)\textasciicircum{}2,lambda(1-lambda),lambda(1-lambda),lambda\textasciicircum{}2). \steptag{validated} \\[4pt]
\step{1.9.4} Because row\_4(A\_lambda)=(lambda,1-lambda), row\_4(A\_lambda M\_lambda)=lambda row\_1(M\_lambda)+(1-lambda)row\_2(M\_lambda)=(lambda(1-lambda),(1-lambda)\textasciicircum{}2,lambda\textasciicircum{}2,lambda(1-lambda)). The four computed rows are all rows of P\_lambda=A\_lambda M\_lambda, proving 1.9. \steptag{validated} \\[6pt]
\step{1.9.4.1} Write u:=row\_1(M\_lambda)=(1-lambda,0,lambda,0) and v:=row\_2(M\_lambda)=(0,1-lambda,0,lambda). By the definition of matrix multiplication and the displayed rows row\_1(A\_lambda)=(1,0), row\_2(A\_lambda)=(0,1), and row\_3(A\_lambda)=(1-lambda,lambda), the first three rows of A\_lambda M\_lambda are respectively 1*u+0*v=u=(1-lambda,0,lambda,0), 0*u+1*v=v=(0,1-lambda,0,lambda), and (1-lambda)u+lambda v=((1-lambda)\textasciicircum{}2,lambda(1-lambda),lambda(1-lambda),lambda\textasciicircum{}2). \steptag{validated} \\[4pt]
\step{1.9.4.2} By the preceding computation, the first three rows have the required forms. Independently, row\_4(A\_lambda)=(lambda,1-lambda), so the definition of matrix multiplication gives row\_4(A\_lambda M\_lambda)=lambda u+(1-lambda)v=(lambda(1-lambda),(1-lambda)\textasciicircum{}2,lambda\textasciicircum{}2,lambda(1-lambda)). These four calculations cover every row because A\_lambda has exactly four rows; hence, with P\_lambda:=A\_lambda M\_lambda, they establish the complete four-row list asserted in node 1.9.4 (and thus its contribution to 1.9). \steptag{validated} \\[4pt]
\step{1.10} All displayed entries of P\_lambda are nonnegative by 1.1; its first two row sums are 1, and each of its last two row sums is (1-lambda)\textasciicircum{}2+2lambda(1-lambda)+lambda\textasciicircum{}2=((1-lambda)+lambda)\textasciicircum{}2=1. Thus P\_lambda is row-stochastic. \steptag{validated} \\[4pt]
\step{1.11} Using 1.9, row\_1(P\_lambda)-row\_3(P\_lambda)=(lambda(1-lambda),-lambda(1-lambda),lambda\textasciicircum{}2,-lambda\textasciicircum{}2). By 1.1 these magnitudes sum to 2lambda(1-lambda)+2lambda\textasciicircum{}2=2lambda, so ||row\_1(P\_lambda)-row\_3(P\_lambda)||\_1=2lambda. \steptag{validated} \\[4pt]
\end{proofbox}

\end{document}
